


\begin{Exercise}[title=Exercice Préparatoire : La Classe CarteBataille]
\begin{Objectif}
    Créer une classe CarteBataille permettant de simuler une confrontation directe entre deux cartes.
Consignes
\end{Objectif}

\hspace{0.5cm}\textbf{I. Définition de la classe et du constructeur}\\ 

\Question  Créez une classe nommée CarteBataille.

\Question Écrivez son constructeur \_\_init\_\_ qui prend deux paramètres : valeur (entier de 2 à 14) et symbole (chaîne de caractères : "Coeur", "Carreau", "Trèfle", "Pique").\\

%
\textbf{II. Méthode de comparaison}\\

\Question Écrivez une méthode battre(self, autre\_carte) qui :
\begin{itemize}
\item Renvoie True si la valeur de la carte actuelle (self) est strictement supérieure à celle de autre\_carte.

    \item Renvoie False sinon.\\
\end{itemize}

%
\noindent\textbf{III. Programme principal (Tests)}\\
Dans votre script principal, réalisez les actions suivantes :\\
\Question Instanciez une carte joueur1 avec un 10 de Trèfle.
\Question Instanciez une carte joueur2 avec un 13 de Cœur (un Roi).
\Question Affichez les deux cartes dans la console.
\Question Utilisez la méthode battre pour vérifier si joueur1 bat joueur2, puis affichez le résultat sous forme d'un message clair (ex: "Joueur 1 gagne !" ou "Joueur 2 gagne !").
\end{Exercise}
